\clearpage
\ubsection{Introduction to the Ordo}{introduction-ordo}

% \vspace{+1em}

	% \ubsubsection{Introduction to the Ordo}{introduction-ordo}
\lettrine[lines=2]{\zallmancaps \color{Red} T}{his} is the companion Ordo to the \emph{Breviarium Predicatorum} of Humbert, or by its full title, \emph{Breviarium cum Psalterio iuxta Ecclesiasticum Officium Sacri Ordinis Fratrum Predicatorum}.

\par This Ordo is laid out in five sections:
\begin{itemize}
\item Celebratory and Navigational Notes
\item The Rubrics, Translated into English
\item The Kalendar for Nov. 2026-2027
\item The Ordo Proper, for Nov. 2026-2027
\item An Appendix with full Specimen Offices
\end{itemize}

\ubsection{Celebratory and Navigational Notes}{celebratory-notes}

	\ubsubsection{On the Office of the Dedication of a Church}{office-dedication}
Please note concerning the Kalendar and Ordo, the Breviary contains rubrics and propers for the Office of the Dedication of a Church, which, for reasons of generality, are placed neither on the Kalendar nor in the Ordo in this book. Those wishing to keep to the rubrics must celebrate the Office of the Dedication when their Church was dedicated, keeping in mind the rules concerning its Octave.

	\ubsubsection{Obituary Announcements}{obit-notes}
In the Ordo, on certain dates, daggers \textdagger \ indicate that at the \emph{Preciosa}, following the rubric after the \emph{Martyrologium}, the Anniversary or Obituary of a Master-General is to be announced for the next day. For example, at the \emph{Preciosa} of February 3rd, it is to be announced that the next day is the Anniversary of Fathers and Mothers, then the \emph{Commemoratio fratrum, etc.}



	\ubsubsection{The \emph{Salve} Procession}{salve-procession}
We read from the \emph{Ordinarium}: ``When the antiphon is begun, all the friars, ... shall bend their knees without prostration, ... remaining on bent knees until [the incipit] \emph{Salve} or \emph{Ave} has been sung, and then, standing upright, they shall go forth in procession to the outer church ... . All brothers, upon going out, shall bow before the cross which is between the choir and the lay church.

While the antiphon is being sung, holy water shall be sprinkled by the hebdomadarian ... . The aforesaid bending of the knees, however, should not be done outside the choir.

When the antiphon is finished, the candle-bearers shall say the versicle: \Vbar . \emph{Dignare me}, adding \emph{Alleluya} in Paschaltide, and let it be said in the manner in which versicles at commemorations are said.

When the versicle is said, the hebdomadarian shall add \emph{Oremus} and the Collect \emph{Concede nos}, in the manner in which Collects at the Hours are said. When the Prayer is finished, let \emph{Fidelium anime} be said.

When this is said, if outsiders are present, the friars shall enter the choir and there say the \emph{Pater noster} and \emph{Credo} in their places with the customary bows or prostrations. If outsiders are not present, they shall say them outside [the choir]. The candle-bearers, however, having always bowed at the saying of \emph{Fidelium}, shall immediately return to the sacristy, saying the \emph{Pater noster} and \emph{Credo} as they go.

If, however, the outer place is not suitable for a procession, or if the friars are notably few in number, it is not necessary to go outside; rather, all things shall be done in the choir, with the candle-bearers nonetheless standing before the steps of the presbytery according to the manner described above.

Where, however, there are outsiders in the church, when any friars have remained behind from Compline, a signal shall be made near the end of Compline so that those who remained behind from Compline may come to the procession.''

In choir, the discipline follows.





	\ubsubsection{On Responsories of First Vespers}{responsories-first-vespers}
At First Vespers of Semiduplexes and greater, as well as Sundays when a History begins, there is a Responsory at First Vespers. Note that the number that follows indicates which lesson it is found after, and that there is always a \emph{Gloria} said with the Responsory, after which the second half from the asterisk (*) is repeated. If a History begins that Sunday, yet it is impeded, it is not reserved for the next Sunday, and is omitted that year.


	\ubsubsection{Notes on the Daily Office of the Blessed Virgin}{daily-office-notes}
At the beginning of the Daily Office of the Blessed Virgin is written again the note that Antiphons are not to be sung partially before the Psalms, but only as a whole after the Psalm(s), as intonation belongs only to the choral office. Similarly, there is a note that if one is not a priest, instead of saying \emph{Dominus vobiscum, etc.}, \emph{Domine exaudi, etc.} replaces it everywhere. Likewise for the \emph{Te Deum} in the Daily Office, it is said when the main office also has a \emph{Te Deum}. Per the 1909 \emph{Breviarium Pr\ae dicatorum}, the third Responsory is repeated in whole, until the versicle exclusive.

As the rubrics and Humbert's Constitution imply, the hours of the Daily Office are said before the main office, except for Compline, when it is said after the main Compline, after the Blessing, but before the \emph{Salve} Antiphon. Such it is also prescribed in the 1909 \emph{Breviarium Pr\ae dicatorum}.


	\ubsubsection{The \emph{Pater noster}}{pater}
The \emph{Pater noster} is written in full only in the \emph{Missale Conventuale} and the \emph{Missale Minorum Altarium}, and not in the preceding books of the Prototype. Here it is per the \emph{Missale Conventuale}.
\begin{multicols}{2}
\fullpsalm{pater}
\end{multicols}


	\ubsubsection{On the Compline Blessing}{compline-blessing}
As friars were never to be absent from Compline, the Blessing for Compline is given by the hebdomadarian and hence not copied in the \emph{Breviarium}. Later Dominican Breviaries include the blessing, and so, in private, the wording is changed from \emph{vos} to \emph{nos}.

	\ubsubsection{The Versicle before Lauds}{lauds versicle}
At Lauds, before the \emph{Deus in adiutorium}, there is daily a versicle that is said, which is either from the Season or from the Sanctoral, excepting the Triduum and All Souls' Day.


	\ubsubsection{On the \emph{Te Deum}}{on-te-deum}
The \emph{Te Deum} is said on all Octaves, Totum Duplex Feasts, and on all Simplex feasts and above when it is neither Septuagesimatide nor Lent. Consequently, on ferial days and three lesson feasts, the \emph{Te Deum} is not said, even in Paschaltide.

	\ubsubsection{On Short Responsories}{short-responsories}
Short Responsories are said thus: The first line, until the versicle exclusive, is said twice, then the versicle, then from the asterisk (*) until the versicle exclusive, then the \emph{Gloria}, without the \emph{Sicut}, then the first line is repeated in full, until the versicle exclusive.

During Passiontide, the \emph{Gloria} is not said, so after the second half from the asterisk (*), the first line is repeated in full.

During Paschaltide and on Totum Duplex feasts, to the Short Responsories are added two \emph{Alleluyas}, which are considered the second half from the asterisk (*), yet, if the feast has an Octave, the Short Responsories are said without the \emph{Alleluya}, if it is outside of Paschaltide.

	\ubsubsection{On Locating Versicles}{locating-versicles}
At times, versicles will be abbreviated or not written in the Breviary. When such things occur, or for similar elements, recourse must be taken to:
\begin{enumerate}[nosep,leftmargin=13pt]
	\item Vespers, Lauds, or the Nocturns at Matins, where they may be written in full,
	\item On some feasts of secondary character, their primary feast: e.g., on Peter in Chains, where some elements are taken from Peter's Chair, and for Assumption of the Blessed Virgin, some things are taken from her Purification.
	\item The Common of Saints, or
	\item For Ferias and Sundays, the first week after the Octave of the Epiphany, or for the Ben. or Mag. Antiphons, the preceding Sunday.
\end{enumerate}

	\ubsubsection{On the Office of the Dead}{office-of-the-dead}
As written in the rubrics, the nine-lesson Office of the Dead is to be said every week, excepting the periods specified and within Octaves. The choice of collects is \emph{ad libitum}, except when prescribed by the Necrology. After the \emph{Requiescant}, the \emph{Pater} is recited, and yet it is not recited before the Office of the Dead, except on All Souls' Day, where it is the main office.

Notably, this is laxer than the rules for the 1909 \emph{Breviarium Pr\ae dicatorum}, which exempts only Totum Duplex feasts and Easter and Pentecost Octaves, and other feasts of precept.

For the three-lesson Office of the Dead, there is no occasion for it in private, as Bonniwell notes that, until the 16th century, when this practice was abolished, it was appointed on a weekly rotation to assigned clerics and hebdomadarian. This was said every day, exempting Saturdays, according to the rubrics. However, Humbert, per Bonniwell, notes that ``Any others who wish to do so, may be present [for the mass and office].''
%that is, Salamanca in 1551


{\color{Red} \ubsubsection{Further Resources on the Dominican Rite}{further-resources}}

For those interested in the history of the Dominican Rite, ``\emph{A History of the Dominican Liturgy}'' by Fr. William Bonniwell, O.P. is the standard. A more recent history from 1945 to 1969 is available online by the Very Rev. Fr. Augustine Thompson, O.P.

Additionally, any questions, typos, or comments can be sent to \emph{BreviaryTypesetting@Outlook.com}



	% \ubsubsection{}{}
	% \ubsubsection{}{}
	% \ubsubsection{}{}
	% \ubsubsection{}{}
	% \ubsubsection{}{}